\section{Enumerations}

\subsection{Section Overview}
Enumerations provide a way to create a user-defined type with a set of named integer constants, called enumerators. This can make code more readable and less error-prone than using raw integers or macros. In this bonus section, we'll cover both unscoped (C-style) enumerations and the modern, safer scoped enumerations (enum classes) introduced in C++11.

\subsection{Enumerations in C++}
An unscoped enumeration is defined with the \texttt{enum} keyword. By default, the enumerators are assigned integer values starting from 0.
\begin{lstlisting}
enum Direction { North, South, East, West };
// North = 0, South = 1, East = 2, West = 3

Direction heading = North;
if (heading == South) {
	// ...
}
\end{lstlisting}
The major drawback of unscoped enums is that their enumerators ``pollute'' the surrounding scope, which can lead to name clashes. They also implicitly convert to integers, which can cause subtle bugs.

\subsection{Coding Exercise 44: Using Enumerations}
Define an unscoped enumeration for \texttt{Color} with enumerators \texttt{Red}, \texttt{Green}, and \texttt{Blue}. Create a variable of this type and assign \texttt{Green} to it.
\begin{lstlisting}
#include <iostream>

enum Color { Red, Green, Blue };

int main() {
	Color my_color = Green;
	if (my_color == Green) {
		std::cout << "The color is green." << std::endl;
	}
	return 0;
}
\end{lstlisting}

\subsection{Scoped Enumerations (enum classes)}
C++11 introduced scoped enumerations (\texttt{enum class}), which solve the problems of unscoped enums.
\begin{itemize}
    \item \textbf{Scoped}: The enumerators are part of the enum's scope, preventing name clashes. You must use the scope resolution operator (e.g., \texttt{Color::Red}).
    \item \textbf{Strongly-typed}: They do not implicitly convert to integers. You must use a \texttt{static\_cast} if you need the integer value.
\end{itemize}
\begin{lstlisting}
enum class Color { Red, Green, Blue };
enum class Traffic_Light { Red, Green, Yellow }; // No name clash with Color::Red

int main() {
	Color my_color = Color::Green;
	// if (my_color == Traffic_Light::Green) // Compile Error! Different types

	int green_value = static_cast<int>(Color::Green); // green_value is 1
}
\end{lstlisting}
You should always prefer \texttt{enum class} over unscoped \texttt{enum} in modern C++.

\subsection{Quiz 19: Section 23 Quiz}
\begin{enumerate}
    \item What is a primary advantage of \texttt{enum class} over unscoped \texttt{enum}?
    \begin{enumerate}
        \item[a)] They can hold floating-point values.
        \item[b)] Their enumerators are scoped and they are strongly-typed, preventing name clashes and unwanted integer conversions.
        \item[c)] They are shorter to type.
        \item[d)] They are compatible with older C code.
    \end{enumerate}

    \item In an unscoped \texttt{enum \{ A, B, C \}}, what is the default integer value of \texttt{B}?
    \begin{enumerate}
        \item[a)] 0
        \item[b)] 1
        \item[c)] 2
        \item[d)] It is not an integer.
    \end{enumerate}

    \item How do you access an enumerator from an \texttt{enum class} named \texttt{Status} with an enumerator \texttt{Pending}?
    \begin{enumerate}
        \item[a)] \texttt{Pending}
        \item[b)] \texttt{Status.Pending}
        \item[c)] \texttt{Status::Pending}
        \item[d)] \texttt{Status->Pending}
    \end{enumerate}
\end{enumerate}

\textbf{Answers:} 1. (b), 2. (b), 3. (c)
